\documentclass{article}
\usepackage{graphicx} % Required for inserting images

\title{Snapshot Objectives}
\author{
Fernando Camberos \\
Kwane Galang \\
Korapat Jaengsrivong \\
Andy Lim \\
Fredy Lung
}
\date{May 2026}

\begin{document}

\maketitle

\section*{Start Objective}

The goal of start objective is to establish the foundation of the Dog Shelter Management System. During this stage, the team will define the core features of the system, organize responsibilities, and outline the overall structure of the project.\\

The main objectives include:
\begin{itemize}
\item Defining core system features, including Dog Intake, Medical Records, Adoption Management, Dog Status Tracking, and an Admin Dashboard
\item Assigning team roles and responsibilities for each component of the project
\item Setting up the GitHub repository and organizing the project structure
\item Creating initial drafts of the SRS, SDD, Design Specification, and User Manual
\item Creating a Jira board and defining initial tasks and sprint structure
\item Identifying required tools, dependencies, and technologies for the system
\end{itemize}

\vspace{2mm}

\noindent At the end of this stage, the project will have a clear direction, defined scope, and structured plan for development.

\vspace{3mm}

\section*{1st Checkpoint Objective}

Building upon the foundation established in the start objective, this stage focused on implementing the first major feature set: the Dog Intake System and Dog Status Tracking functionality. The objective of this checkpoint was to establish the system’s ability to manage and track dogs within the shelter.\\

The main objectives include:
\begin{itemize}
\item Implementing the Dog Intake System, allowing staff to add and store basic dog information
\item Implementing Dog Status Tracking (Available, Pending, Adopted)
\item Creating and assigning new Jira tasks related to these features
\item Developing TestRail test cases to validate dog intake and status updates
\item Updating the SRS to include detailed requirements for dog intake and tracking
\item Updating the SDD to describe how these features are structured and managed
\item Expanding the Design Specification to include dog-related pages and interactions
\item Updating the User Manual to reflect new functionality
\item Updating the workflow diagram to reflect how dog data flows through the system
\end{itemize}

\vspace{2mm}

\noindent At this stage, the system establishes its core functionality for managing dogs within the shelter.

\vspace{3mm}

\section*{2nd Checkpoint Objective}

Following the successful implementation of the initial system features, this stage expanded the system to include additional functionality.\\

The main objectives include:
\begin{itemize}
\item Implementing the Medical Records System, including vaccination tracking, medical history, and updating health status
\item Implementing the Adoption Management System, allowing users to submit applications and staff to review them
\item Implementing the Admin Dashboard to allow staff to manage dogs, records, and applications in a centralized interface
\item Creating and assigning new Jira tasks related to medical, adoption, and dashboard features
\item Developing additional TestRail test cases for medical records, adoption workflows, and dashboard functionality
\item Updating the SRS with requirements for medical records, adoption management, and admin dashboard features
\item Updating the SDD to include design details for these systems and dashboard interactions
\item Expanding the Design Specification to include medical record pages, adoption workflows, and admin dashboard interfaces
\item Updating the User Manual to include instructions for the newly added features
\item Updating the workflow diagram to reflect new data interactions and system components
\end{itemize}

\vspace{2mm}

\noindent At the end of this stage, the system supports both operational tracking and user interaction through adoption workflows.

\vspace{3mm}

\section*{Final / Due Date Checkpoint}

After completing the core features of the system, the final stage focused on refinement, testing, and identifying future improvements.\\

The main objectives include:
\begin{itemize}
\item Refining all implemented features and ensuring consistency across the system
\item Fixing any identified issues and improving usability
\item Creating final TestRail reports and summarizing testing results
\item Updating all documents (SRS, SDD, Design Specification, User Manual) for completeness and accuracy
\item Finalizing the workflow diagram to reflect the complete system
\item Documenting future work, including potential features that were not implemented, such as enhanced admin tools or notification systems
\end{itemize}

\vspace{2mm}

\noindent At the end of this stage, the project was considered complete, with all major features documented and a clear outline of future enhancements that could further improve the system.

\end{document}